\chapter{Developed Solution}\label{chap:chap3}

% Este capítulo deve começar por fazer uma apresentação detalhada do problema a resolver podendo mesmo, 
% caso se justifique, constituir-se um capítulo separado só com essa finalidade.

% Defining the problem: BT, Spectral and Score aligned visualization.
% Defining the proposed solution: A visualization tool that can present BT and Audio data with an aligned musical score.

% Developed solution: Tool for generating custom visualizations that present beat data, spectrogram and score all aligned. 

% System Architecture: Python based tool, that employs a modular and Object Oriented architecture for generating the visualizations.

% Choices made and why
\section{MIRVisualScore}
MIRVisualScore is intended to be a Python OO tool for visualizing recorded performances synchronized with the corresponding music score in a time-aligned manner.

\section{Developed Iterations}
This thesis started with the proposed title ``Interactive Musical Beat Annotation GUI''. The initial thesis proposition was to evaluate modern GUI software solutions used for beat annotation (such as  SV), and build upon them with a focus in beat annotation practices.
The proposed goals included:
\begin{enumerate}
    \item Either develop a web-based annotation system, or advancing the existing SV framework. Emphasizing collaborative features, real-time feedback and build on interaction mechanisms for beat annotation. 
    \item Integration of ML human-in-the-loop features that provide intelligent annotation suggestions. Inspired by 
\end{enumerate}

% ==============================================================================
% INITIAL THESIS PROPOSITION
% ``Interactive Musical Beat Annotation GUI'' 
% ## Goals
% 1. Modernize Annotation Infrastructure: Develop a contemporary web-based annotation system or advance the existing Sonic Visualizer framework, incorporating collaborative features and real-time feedback and interaction mechanisms that surpass current standalone tool limitations.
% 2. As part of the system architecture phase, a Python wire-frame library will be shared early in the project to visualize and test annotation concepts, enabling community feedback before final interface decisions.
% 3. Integrate Self-Supervised Learning: Implement self-supervised learning (SSL) techniques that provide intelligent annotation suggestions without requiring extensive pre-labeled datasets, addressing the cold-start problem effectively.
% 4. Create Interactive ML Feedback Loops: Design user-centric iteration mechanisms where user annotations inform model refinement in real-time, creating a symbiotic human-machine annotation process.
% 5. Validate Through User Studies: Conduct comprehensive testing with diverse users (musicologists, producers, researchers) to assess effectiveness, usability, and accuracy improvements across various musical scenarios.
% 6. Open Source Impact: Contribute to the research community through GitHub repositories and libraries, potentially integrating with existing frameworks like Gradio or Hugging Face, and target publication at conferences like Intelligent User Interfaces (IUI), ACM Multimedia or International Society for Music Information Retrieval (ISMIR).

% ## Innovative Aspects
% 1. Enhanced Annotation Experience: Move beyond traditional desktop-bound annotation workflows by creating an integrated system that actively engages users in semi-automatic annotation processes, leveraging modern technologies for improved collaboration and accessibility. 
% 2. Intelligent Visual Feedback: Develop interactive interfaces where signal processing techniques (e.g. tempo grid alignment, onset detection) and ML model outputs (e.g. beat activation functions) are dynamically visualized alongside user annotations, providing immediate visual correlation and conflict detection.
% 3. Bidirectional Learning: Create a system where user expertise refines ML models in real-time while ML predictions inform and accelerate user annotation decisions, establishing an effective human-machine collaboration workflow. 
% 4. Cold-Start Mitigation: Implement streamlined workflows for annotating challenging musical excerpts, using SSL techniques to bootstrap model performance from minimal initial annotations.
% ==============================================================================

\subsection{BeatMapJS}
% Goal: 
% Tools Used:
% Developed Features: 
% Conclusion:

\subsection{MEI-Basic-Edit}
% Goal: 
% Tools Used:
% Developed Features: 
% Conclusion:
\subsection{Score-Allignment.py}
% Goal: 
% Tools Used:
% Developed Features: 
% Conclusion:
\subsection{Score-Allignment.js}
% Goal: 
% Tools Used:
% Developed Features: 
% Conclusion: